% !PRINTED 407
\weberpage{407}
Every part of the sum on the left side of formula (5) which
refers to an ideal class then gives the same contribution
to the sum, and we obtain:

\weq{7}{\begin{aligned}
\operatorname{Lim}_{s=1}(s-1)\sum\frac{1}{m^s}\sum\frac{(\Delta,n)}{n^s} &= h\,\frac{1}{\sqrt{\Delta}}\log\varepsilon, & \Delta&>0\\
&= h\,\frac{\pi}{\sqrt{-\Delta}}, & \Delta&<0.
\end{aligned}}

\gesperrt{Here $h$ is the class number of the field $\Omega$, and in order
not to have to cite the two exceptional cases $\Delta = -4$,
$\Delta = -3$ again and again, we shall stipulate that
in these two cases, in which the class number is 1,
the value $\frac{1}{2}$ or $\frac{1}{3}$ is to be understood by $h$.}

As often as $n$ runs through a complete system of residues to the modulus $\Delta$
, we have $\Sigma(\Delta,n) = 0$ [§ 85, (22)], and consequently the sum
$\Sigma(\Delta,n)$, when the $n$ are arranged in order of magnitude,
always remains finite, however far $n$ may increase. Accordingly the theorem Bd.\,II,
§ 196, 1. is applicable, according to which

\[
\operatorname{Lim}_{s=1}\sum\frac{(\Delta,n)}{n^s} = \sum\frac{(\Delta,n)}{n}
\]

holds. On the other hand

\[
\operatorname{Lim}(s-1)\Sigma\frac{1}{m^s} = 1
\]

(Bd.\,II, § 196, p. 723), and from (7) there results:

\weq{8}{\begin{aligned}
h\log\varepsilon &= \sqrt{\Delta}\sum_{1,\infty}^{n}\frac{(\Delta,n)}{n}, & \Delta&>0\\
h\pi &= \sqrt{-\Delta}\sum_{1,\infty}^{n}\frac{(\Delta,n)}{n}, & \Delta&<0.
\end{aligned}}

The sum

\[
\sigma = \Sigma\frac{(\Delta,n)}{n}
\]

can be represented in finite form. For, by
§ 86, (2):

\weq{9}{\sqrt{\Delta}\,(\Delta,n) = \overset{k}{\Sigma}(\Delta,k)\,e^{-\frac{2nk\pi i}{\Delta}},}

in which $k$ runs through a complete system of residues to the modulus $\Delta$
 and, for a negative discriminant, $\sqrt{\Delta} = i\sqrt{-\Delta}$ is to be
set. Accordingly it follows:

\[
\sigma = \frac{1}{\sqrt{\Delta}}\sum_{}^{k}(\Delta,k)\sum_{}^{n}\frac{1}{n}\,e^{-\frac{2nk\pi i}{\Delta}},
\]
%CONT>>
